    %   Il progetto nasce dal template per il frontespizio realizzato da Marco Antonio Corallo, che ringrazio. Seguono alcuni commenti per evidenziare la presenza di alcuni pacchetti che mi sono stati utili per la stesura della tesi. Chiaramente, dipende tutto dal tipo di lavoro che uno vuole eseguire, che determina anche le diverse esigenze. Durante la stesura ho passato molto tempo su siti e forum a cercare di risolvere alcuni probelmi di formattazione, ma in generale Latex è stato piuttosto versatile. 

% Tipo di documento. L'uso di twoside implica che i capitoli inizino sempre con la prima pagina a sinistra, eventualmente lasciando una pagina vuota nel capitolo precedente. Se questa cosa è fastidiosa, è possibile rimuoverlo. 
\documentclass[a4paper,openright]{report}

% Dimensione dei margini
\usepackage[a4paper,top=3cm,bottom=3cm,left=3cm,right=3cm]{geometry} 
% Dimensione del font
\usepackage[fontsize=13pt]{scrextend}
% Lingua del testo
\usepackage[english,italian]{babel}
% Lingua per la bibliografia
\usepackage[fixlanguage]{babelbib}
% Codifica del testo
\usepackage[utf8]{inputenc} 
% Encoding del testo
\usepackage[T1]{fontenc}
% Permette di generare testo fittizio. Mi è stato utile 
% per capire quale sarebbe stata l'impostazione del 
% testo nella pagina prima che scrivessi un determinato paragrafo
\usepackage{lipsum}
% Per ruotare le immagini
\usepackage{rotating}
% Per modificare l'header delle pagine 
\usepackage{fancyhdr}               

% Librerie matematiche
\usepackage{amssymb}
\usepackage{amsmath}
\usepackage{amsthm}         

% Uso delle immagini
\usepackage{graphicx}
% Uso dei colori
\usepackage[dvipsnames]{xcolor}         
% Uso dei listing per il codice
\usepackage{listings}          
% Per inserire gli hyperlinks tra i vari elementi del testo 
\usepackage{hyperref}     
% Diversi tipi di sottolineature
\usepackage[normalem]{ulem}

% -----------------------------------------------------------------

% Modifica lo stile dell'header
\pagestyle{fancy}
\fancyhf{}
\lhead{\rightmark}
\rhead{\textbf{\thepage}}
\fancyfoot{}
\setlength{\headheight}{12.5pt}

% Rimuove il numero di pagina all'inizio dei capitoli
\fancypagestyle{plain}{
  \fancyfoot{}
  \fancyhead{}
  \renewcommand{\headrulewidth}{0pt}
}

% Stile del codice
\lstdefinestyle{codeStyle}{
    % Colore dei commenti
    commentstyle=\color{teal},
    % Colore delle keyword
    keywordstyle=\color{Magenta},
    % Stile dei numeri di riga
    numberstyle=\tiny\color{gray},
    % Colore delle stringhe
    stringstyle=\color{violet},
    % Dimensione e stile del testo
    basicstyle=\ttfamily\footnotesize,
    % newline solo ai whitespaces
    breakatwhitespace=false,     
    % newline si/no
    breaklines=true,                 
    % Posizione della caption, top/bottom 
    captionpos=b,                    
    % Mantiene gli spazi nel codice, utile per l'indentazione
    keepspaces=true,                 
    % Dove visualizzare i numeri di linea
    numbers=left,                    
    % Distanza tra i numeri di linea
    numbersep=5pt,                  
    % Mostra gli spazi bianchi o meno
    showspaces=false,                
    % Mostra gli spazi bianchi nelle stringhe
    showstringspaces=false,
    % Mostra i tab
    showtabs=false,
    % Dimensione dei tab
    tabsize=2
} \lstset{style=codeStyle}

% Stile di codice per dimensioni maggiori, in cui ho avuto bisogno di un testo più picolo (ad esempio se si vuole inserire del codice che ha linee molto lunghe). Per usare questo stile piuttosto che il precedente, usare 

% \lstset{style=longBlock}
%  % inserire il codice...
% \lstset{style=codeStyle}

% Il secondo comando consente di tornare allo stile precedente 
\lstdefinestyle{longBlock}{
    commentstyle=\color{teal},
    keywordstyle=\color{Magenta},
    numberstyle=\tiny\color{gray},
    stringstyle=\color{violet},
    basicstyle=\ttfamily\scriptsize,
    breakatwhitespace=false,         
    breaklines=true,                 
    captionpos=b,                    
    keepspaces=true,                 
    numbers=left,                    
    numbersep=5pt,                  
    showspaces=false,                
    showstringspaces=false,
    showtabs=false,                  
    tabsize=2
} \lstset{style=codeStyle}

% Togliendo il commento al comando che segue, si inseriscono nella bibliografia anche le fonti presenti in Bibliography.bib ma non citati direttamente con il comando \cite
% \nocite{*}

% Margini prima e dopo blocchi di codice, per avere più distanza
\lstset{aboveskip=20pt,belowskip=20pt}

% Modifica dello stile dei riferimenti, con il testo in cyano
\hypersetup{
    colorlinks,
    linkcolor=CornflowerBlue,
    citecolor=CornflowerBlue
}

% Aggiunti definizioni, teoremi, linea e listing
\newtheorem{definition}{Definizione}[section]
\newtheorem{theorem}{Teorema}[section]
\providecommand*\definitionautorefname{Definizione}
\providecommand*\theoremautorefname{Teorema}
\providecommand*{\listingautorefname}{Listing}
\providecommand*\lstnumberautorefname{Linea}

\raggedbottom

%\newcommand{\cgs}[1]{{\textcolor{brown}[\textcolor{red}{\bf{GS: }}{ \textcolor{brown}{#1]}}}}             
%\newcommand{\cmc}[1]{{\textcolor{blue}[\textcolor{magenta}{\bf{MC: }}{ \textcolor{blue}{#1]}}}}



% -----------------------------------------------------------------
\newread\tmp

\newcommand{\quickcharcount}[1]{%
  \immediate\write18{texcount -1 -sum -merge -char #1.tex > #1-chars.sum}%
  \openin\tmp=#1-chars.sum%
  \read\tmp to \thechar%
  \closein\tmp%
}

\newcommand{\quickwordcount}[1]{%
  \immediate\write18{texcount -1 -sum -merge #1.tex > #1-words.sum}%
  \openin\tmp=#1-words.sum%
  \read\tmp to \theword%
  \closein\tmp%
}



\begin{document}

\begin{titlepage}
\begin{figure}[!htb]
    \centering
    \includegraphics[keepaspectratio=true,scale=0.5]{images/Frontespizio/cherubinFrontespizio.eps}
\end{figure}

\begin{center}
    \LARGE{UNIVERSITÀ DI PISA}
    \vspace{5mm}
    \\ \large{DIPARTIMENTO DI FISICA "ENRICO FERMI" }
    \vspace{5mm}
    \\ \LARGE{Laurea Triennale In Fisica}
\end{center}

\vspace{15mm}
\begin{center}
    {\LARGE{\bf Impiego di transmoni superconduttori \\ \vspace{1mm} per la rivelazione \\ \vspace{3mm} di    assioni e materia oscura
 }}
    
    % Se il titolo è abbastanza corto da stare su una riga, si può usare
    
    % {\LARGE{\bf Un fantastico titolo per la mia tesi!}}
\end{center}
\vspace{30mm}

\begin{minipage}[t]{0.47\textwidth}
	{\large{Relatore:}{\normalsize\vspace{3mm}
	\bf\\ \large{Prof: Gianluca Lamanna}}}
\end{minipage}
\hfill
\begin{minipage}[t]{0.47\textwidth}\raggedleft
	{\large{Candidato:}{\normalsize\vspace{3mm} \bf\\ \large{Lorenzo Zaffina}}}
\end{minipage}

\vspace{30mm}
\hrulefill
\\\centering{\large{Anno Accademico 2021/2022}}

\end{titlepage}
\section*{Long Abstract}
Uno dei problemi aperti della fisica moderna è quello riguardante la materia oscura. A supporto della sua esistenza vi sono innumerevoli osservazioni astronomiche, ma la sua natura fondamentale rimane ignota. Una delle più promettenti proposte per la spiegazione di tale fenomeno è rappresentata dalla possibile esistenza di ipotetiche particelle chiamate assioni.\\ %\cite{experimentalsearches} \cite{hunt}\\

L’idea degli assioni venne introdotta nel 1977 da Peccei e Quinn per risolvere il problema della mancata osservazione della violazione di CP nell'interazione forte tra quark. Queste particelle (insieme alle ALPs, Axion-Like Particles) dovrebbero interagire estremamente poco con le particelle del Modello Standard e perciò rappresentano dei validi candidati per la materia oscura.\\ 


Per la rivelazione degli assioni si sfrutta l’accoppiamento del campo assionico con il campo elettromagnetico (EM). Gli sforzi sperimentali vengono orientati sia alla ricerca di assioni di origine cosmica, sia in esperimenti condotti puramente in laboratorio.

%Per la rivelazione degli assioni si sfrutta l’accoppiamento del campo assionico con il campo Elettromagnetico (EM) e le ricerche sperimentali vengono essenzialmente condotte su due fronti: da un lato la ricerca di assioni di origine cosmica, dall'altro esperimenti condotti puramente in laboratorio.\\ 


%In particolare, la costante di accoppiamento $g_{a\gamma\gamma}$ è proporzionale alla massa a riposo dell’assione $m_ac^2$.
%Nel primo caso l’idea è quella di rivelare i fotoni prodotti dalla riconversione degli assioni generati da fonti astronomiche. 
Ad esempio, gli elioscopi servono a rivelare assioni prodotti nel nucleo solare, riconvertendoli in fotoni attraverso potenti magneti puntati verso il sole. Gli aloscopi invece hanno l'obiettivo di rivelare gli assioni che dovrebbero popolare l’alone galattico.\\
I principali esperimenti in laboratorio sono i “Light Shining through a Wall” (LSW), basati sulla generazione di assioni e sulla rigenerazione di fotoni tramite intensi campi magnetici. \\

In tutti questi casi, la rivelazione di assioni si basa sull’accoppiamento con i fotoni, tipicamente nelle microonde.  %Siccome tali particelle sono estremamente poco interagenti, la sfida sperimentale è quella di riuscire a rivelare singoli fotoni, .\\
Tali fotoni, nel caso di  LSW e aloscopi, vengono accumulati in cavità risonanti con frequenza di risonanza opportuna. Nelle condizioni sperimentali tipiche si stima un numero medio di fotoni pari a $\Bar{n}_{axion} \sim 10^{-8}-10^{-5}$ per misurazione \cite{PhysRevLett.126.141302}. La sfida sperimentale è dunque quella di riuscire a rivelare anche singoli fotoni nelle microonde.\\ Attualmente questi esperimenti utilizzano amplificatori lineari che operano al limite dello SQL (Standard Quantum Limit), le cui fluttuazioni corrispondono ad un background effettivo pari a $\Bar{n}_{SQL}=1$. 
Il rumore quantistico copre completamente il segnale ($\Bar{n}_{SQL}\gg\Bar{n}_{axion} $), rendendo impossibile la rivelazione.\\

Anziché utilizzare amplificatori lineari, si vogliono sviluppare dei rivelatori sensibili ad un singolo fotone. Esistono già delle tecnologie adatte a rivelare singoli fotoni nell’infrarosso, ma non sono altrettanto adatte a misurare fotoni di bassa energia nelle microonde.\\% come nel nostro caso.\\

È qui che entrano in gioco i transmoni, un tipo di qubit superconduttori costituiti da una giunzione Josephson e un condensatore di shunt.\\
Tali dispositivi sono alla base dei computer quantistici, ma giocano un ruolo chiave anche nella rivelazione di fotoni, sfruttando l’interazione
tra il qubit e il campo EM in una cavità risonante.

%sfruttando l'accoppiamento del con il campo EM in una cavità risonante.
%per rivelare i singoli fotoni si sfrutta l’interazione tra un qubit superconduttore e il campo EM nella cavità risonante. 
Infatti, la frequenza di transizione del qubit è legata al numero di fotoni nella cavità. Questo permette di sviluppare un protocollo di misura che possa inferire il numero di fotoni nella cavità a partire da una misura dello stato del qubit.\\

Il punto di forza di questo metodo è che si tratta di una misurazione "quantum non-demolition" ovvero tale da non perturbare il sistema quantistico che va ad indagare. Questo permette di superare lo SQL, dando un impulso notevole alla ricerca di particelle assioniche.% incrementando la velocità di ricerca di assioni di più ordini di grandezza.





\quickwordcount{chapters/LongAbstract}
\quickcharcount{chapters/LongAbstract}

%There are \thechar characters and approximately \theword spaces.
%That makes approximately \the\numexpr\theword+\thechar\relax\ characters total.


%Numero di caratteri: \the\numexpr\theword+\thechar\relax\ /3000 $\pm$ 300

%
\section*{Bibliografia prova}


[1] Graham et al., Annu. Rev. Nucl. Part. Sci. 2015. 65:485–514; Experimental Searches for the Axion and Axion-Like Particles.\\

\noindent
[2] Andreas Ringwald, arXiv:1506.04259v1  ; The hunt for axions




\appendix
\bibliographystyle{plain}
\bibliography{chapters/Bibliografia.bib}
\nocite{*}
\end{document}
% -----------------------------------------------------------------
